\chapter[Introduction]{Introduction}\label{chapter:introduction}

\begin{quote}\itshape
The purpose of this chapter is to introduce the research problem addressed by
this thesis and explain how the study investigates it. It begins by motivating
the need to synthesise robot-specific drivers that implement reusable software
interfaces through which high-level controllers can operate different robots,
rather than assuming that executable robot integrations already exist. The
chapter then defines the research objectives and questions, outlines the three
experiments used to answer them, and states the scientific contributions
supported by the study.
\end{quote}

\section[Motivation]{Motivation}
\label{sec:intro-motivation}

Recent robot-learning systems have made task-level behaviour more reusable.
Visuomotor policies learn actions from observations
\citep{chi2023diffusionpolicy}, while vision--language--action models seek to
transfer behaviour across tasks and robots
\citep{brohan2023rt2,openxembodiment2024}. Large language model (LLM) systems
can also interpret instructions and compose supplied robot operations
\citep{ahn2022saycan,liang2023codeaspolicies}. These advances reduce
task-specific programming, but they generally assume that an executable
software interface to the target robot already exists.

The executable implementation behind that interface must translate requested
operations into the joint commands, actuator inputs, coordinate frames,
limits, and observations used by a particular robot. These details vary with
the robot and its control configuration, so a new platform can require new
integration even when the task-level logic is unchanged. Robotics frameworks
expose this remaining work directly: adding hardware still requires
connections, calibration, state reading, and command writing
\citep{lerobot2026hardware,ros2control2026hardware}. This thesis calls the
decision-making program that chooses robot operations the
\emph{high-level controller}, and the reusable interface that hides
robot-specific execution details the \emph{robot capability layer}. The
executable implementation of that layer for one declared robot configuration
is the \emph{robot-specific driver}. Each operation exposed by the layer is a
\emph{capability}; the word \emph{skill} is reserved for a task-level,
temporally extended behaviour such as a pick-and-place sequence.

AutoAdapter 1.0 provided initial evidence that an LLM can generate a
robot-specific driver for this missing interface from machine-readable robot
descriptions. Across eight simulated robot configurations, it produced
self-contained drivers that all passed its
structural-and-smoke onboarding harness, at a mean LLM-inference cost of
US\$\(2.98 \pm 0.84\) and \(7.8 \pm 1.6\) minutes of wall-clock time per robot
\citep{wang2026autoadapter}. This onboarding result did not establish full
behavioural assurance. The same study also separated capability-layer use from
driver synthesis: all seven evaluated LLMs could operate through the same
fixed capability layer backed by the same fixed and validated driver, whereas
only three produced a driver that passed driver validation
\citep{wang2026autoadapter}. Using an available robot capability layer therefore
does not establish that the same model can synthesise a driver that implements
it.

Auto-Adapter 2.0 addresses decisions that the earlier system left
prespecified. Instead of receiving a pre-authored capability catalogue, the
model must derive five to ten reusable capabilities from at least twenty
source-backed tasks with reviewed task-level pass standards. It must then
design measurable conditions for accepting each capability---the
\emph{capability-level pass criteria}---without weakening those standards.
Robot code can run while controlling the wrong joints, using incorrect units,
or producing the wrong physical outcome: a form of the software test-oracle
problem \citep{barr2015oracle,araujo2023rasreview}. Auto-Adapter therefore uses
an independent evaluation program, the \emph{Harness}, to own the private
validation cases, observe simulator outcomes, and issue final verdicts.
Structured failure reports may guide bounded corrections to the current
driver, called \emph{Repair}; reviewed evidence from completed runs may instead inform
later runs through \emph{Evolution}, but cannot alter the run that produced
it.

Three uncertainties then become central. First, LLM backbones may differ both
in synthesising a robot-specific driver that implements a capability layer and
in using the same fixed layer backed by the same fixed and validated driver
through the same high-level-controller implementation. Second, the value of
task-grounded capability design, source-grounded capability-level pass-criteria
design, and continued Evolution must be established rather than assumed.
Third, different physical forms and joint arrangements---their
\emph{morphologies}---may present different difficulties when the model
synthesises robot-specific drivers that implement operations producing robot
motion. This thesis terms these operations \emph{low-level motion-control
capabilities}. With the available robot cases, this comparison can identify
associations and failure patterns, but cannot establish a causal morphology
effect.

These uncertainties motivate a controlled study of when Auto-Adapter can
design reusable robot capability layers and generate and validate their
robot-specific drivers. The thesis evaluates driver synthesis and
capability-layer use separately, analyses the framework's principal
components, and compares driver synthesis across robot morphologies under
matched conditions. The next section converts this motivation into explicit
research objectives and questions.


\section[Research Objectives]{Research Objectives}
\label{sec:research-objectives}
\label{sec:research-questions}

The aim of this thesis is to determine, under controlled conditions, when LLMs
can synthesise robot-specific drivers and use reusable robot capability layers.
Auto-Adapter provides the experimental framework for examining three sources
of variation: the LLM backbone, selected framework mechanisms, and robot
morphology. Chapter~3 uses AutoAdapter-Bench to evaluate robot-specific driver
synthesis and capability-layer use as separate outcomes. Evidence that a model
can operate an available capability layer does not establish that it can
synthesise a driver implementing that layer, or vice versa.

This aim is expressed through three research questions:

\begin{enumerate}
    \item \textbf{RQ1---Comparison of LLM backbones in robot-specific driver
    synthesis and capability-layer use.}
    To what extent do LLM backbones differ in synthesising a robot-specific
    driver that implements a reusable robot capability layer under
    skeleton-assisted and from-scratch generation conditions, and in using the
    same fixed capability layer backed by the same fixed and validated driver
    through the same high-level-controller implementation to complete matched
    tasks?

    \item \textbf{RQ2---Component analysis of the Auto-Adapter framework.}
    Do selected Auto-Adapter mechanisms produce measurable differences in
    robot-software synthesis quality?
    \textbf{(a)} Can an LLM derive five to ten reusable capabilities from at
    least twenty source-backed tasks without receiving a pre-authored
    capability catalogue?
    \textbf{(b)} Can it formulate source-grounded capability-level pass
    criteria, without human predefinition at the capability level, that an
    independent and implementation-blind process can audit and compile into
    executable validation?
    \textbf{(c)} Does reviewed evidence from completed runs improve
    robot-software synthesis quality in matched later runs relative to the same
    condition without that Experience?

    \item \textbf{RQ3---Cross-morphology analysis of robot-specific driver
    synthesis for low-level motion-control capabilities.}
    With the LLM backbone, high-level-controller implementation, Auto-Adapter
    configuration, generation condition, evaluation protocol, and resource
    budget held fixed, what associations are observed between robot morphology
    and the validation success, failure patterns, and resource use of
    synthesised robot-specific drivers that implement low-level motion-control
    capabilities?
\end{enumerate}

\label{sec:intro-scope}
These questions concern declared robot configurations evaluated through the
Direct-MuJoCo simulation route. The resulting evidence can support controlled
comparisons and cross-morphology associations within this scope, but not
hardware validity, sim-to-real transfer, physical validation of unselected
tasks, universal model or robot superiority, or a causal morphology effect.
The following section translates these questions into three controlled
experiments.


\section[Research Experiments]{Research Experiments}
\label{sec:research-experiments}
\label{sec:research-programme}

The thesis begins with a broad comparison of LLM backbones in
robot-specific driver synthesis and capability-layer use through
AutoAdapter-Bench. It then narrows the analysis to three mechanisms within
Auto-Adapter: task-grounded capability design, source-grounded
capability-level pass-criteria design, and Continued Evolution. Finally, it
extends the driver-synthesis analysis across robot morphologies under matched
conditions. The empirical programme is organised into three experiments,
reported in Chapters~3--5:

\begin{enumerate}[leftmargin=*,itemsep=1em]
    \item \textbf{Experiment 1---Comparison of LLM backbones in
    robot-specific driver synthesis and capability-layer use
    (AutoAdapter-Bench, Chapter~3).}
    AutoAdapter-Bench compares the same set of LLM backbones in two tracks.
    The \textbf{driver-synthesis track} holds the Auto-Adapter framework,
    robot inputs, source-backed Task Library, source-task pass standards,
    private-case construction policy, Harness measurement and verdict rules,
    and resource budget fixed within each generation condition. For each
    backbone and robot configuration, one robot capability layer design is
    sealed and used in both generation conditions; the backbone then
    synthesises a corresponding robot-specific driver in each condition. The
    conditions differ only in authorised access to the trusted skeleton. We
    compare independent driver-validation outcomes, failure patterns, and
    resource use; the capability-level pass criteria authored by each model
    remain evaluated outputs rather than fixed inputs. The \textbf{use track}
    gives the same backbones a fixed capability layer backed by the same fixed
    and validated robot-specific driver, the same high-level-controller
    implementation, matched tasks, evaluation protocol, and budget. It compares
    task completion while preventing successful capability-layer use from
    being interpreted as driver-synthesis evidence.

    \item \textbf{Experiment 2---Component analysis of the Auto-Adapter
    framework (Chapter~4).}
    This experiment examines three framework mechanisms and the evidence each
    must produce.
    \textbf{(a) Task-grounded capability design.} We test whether an LLM can
    derive five to ten reusable capabilities from at least twenty
    source-backed tasks without receiving a pre-authored capability catalogue
    or task-to-capability mapping. The analysis measures coverage of the
    source-backed tasks and assesses whether the grouping represents reusable
    robot operations rather than task-specific programs.
    \textbf{(b) Source-grounded capability-level pass-criteria design.} The
    Task Library supplies human-reviewed task-level pass standards derived from
    traceable primary robotics benchmarks or industrial-production standards.
    We test whether an LLM can translate these obligations into
    source-preserving criteria for the capabilities it designs, without
    receiving human-authored capability-level criteria. An independent,
    implementation-blind process then audits the criteria against their source
    obligations and compiles the accepted criteria into executable validation
    cases.
    \textbf{(c) Continued Evolution.} We compare matched later runs with and
    without reviewed Experience derived from completed runs, and measure
    whether access to that Experience improves robot-software synthesis quality
    in later runs.

    \item \textbf{Experiment 3---Cross-morphology analysis of robot-specific
    driver synthesis for low-level motion-control capabilities (Chapter~5).}
    We compare the synthesis of robot-specific drivers that implement
    low-level motion-control capabilities across declared robot configurations
    with different morphologies, while holding the LLM backbone,
    high-level-controller implementation, Auto-Adapter configuration,
    generation condition, evaluation protocol, and resource budget fixed.
    Validation success, failure patterns, and resource use are compared across
    morphologies. This design identifies descriptive associations between
    morphology and driver-synthesis outcomes; it does not establish that
    morphology causes those differences.
\end{enumerate}

Together, the experiments provide the evidence used to answer the three
research questions. The following section states the scientific contributions
that this evidence is designed to support.


\section[Scientific Contributions]{Scientific Contributions}
\label{sec:intro-contributions}

This thesis makes three scientific contributions and three bounded,
industry-relevant contributions. The scientific contributions state what the
research adds to the study of LLM agents and robot-software synthesis. The
industry-relevant contributions state what engineering teams can use within
the declared Direct-MuJoCo scope.

The scientific contributions are as follows:

\begin{enumerate}[leftmargin=*,itemsep=0.8em]

  \item AutoAdapter-Bench: separating robot-specific driver synthesis from
  capability-layer use.

  Tool-use research typically measures whether a model can call an interface
  that is already available
  \citep{ahn2022saycan,liang2023codeaspolicies,huang2023voxposer}.
  AutoAdapter-Bench asks a different question: can the same model synthesise a
  robot-specific driver implementing that interface? It compares the same LLM
  backbones in separate tracks for driver synthesis and capability-layer use.
  This makes robot-specific driver synthesis a separate empirical object and
  prevents successful use from being counted as driver-synthesis evidence; it
  does not assume that the two outcomes are statistically independent.

  \item Auto-Adapter: making capability and pass-criteria design explicit
  components of robot-software synthesis.

  Existing robot-code systems usually assume that the available robot
  operations, the tests used to judge them, or both have already been defined
  by humans
  \citep{liang2023codeaspolicies,tripathi2026testdriven,
  kim2026physicalcorrection}.
  Auto-Adapter removes both assumptions in the setting studied here. Instead
  of receiving a capability catalogue or human-authored capability-level pass
  criteria, the LLM derives reusable capabilities from at least twenty
  source-backed tasks and derives the criteria from reviewed task-level
  standards. An implementation-blind process audits and compiles those
  criteria, and the independent Harness---not the model that wrote the
  driver---controls validation. Auto-Adapter also separates Repair, which can
  change the current driver, from Evolution, which carries reviewed evidence
  into later runs. This makes capability design, pass-criteria design, driver
  synthesis, Repair, and Evolution separate research objects rather than one
  opaque code-generation step. Successful compilation alone is not treated as
  evidence that a criterion is reliable.

  \item Cross-morphology analysis of robot-specific driver synthesis for
  low-level motion-control capabilities.

  Cross-embodiment research usually asks whether a policy can act through
  supplied interfaces on different robots
  \citep{openxembodiment2024,octo2024}.
  This thesis asks a lower-level question: how does synthesising a
  robot-specific driver for low-level motion-control capabilities differ across
  robot morphologies? It holds the LLM backbone, high-level-controller
  implementation, Auto-Adapter configuration, generation condition, evaluation
  protocol, and resource budget fixed, then compares validation outcomes,
  failure patterns, and resource use. This shifts the analysis from using an
  existing interface to synthesising the robot-specific driver that implements
  it. Because morphology co-varies with actuation, dynamics, task applicability,
  and MuJoCo control structure, the results are interpreted as
  associations rather than causal effects.

\end{enumerate}

Within the declared Direct-MuJoCo scope, the bounded industry-relevant
contributions are as follows:

\begin{enumerate}[leftmargin=*,itemsep=0.8em]

  \item From reviewed task standards to executable capability validation.

  Auto-Adapter gives robot testing teams a direct path from documented task
  requirements to executable simulator validation. It records how each
  human-reviewed task-level pass standard is translated into model-authored
  capability-level pass criteria, how those criteria are independently audited
  and compiled, and how the resulting validation cases are evaluated by the
  Harness. The resulting source links, audit records, and
  verdicts make the basis of each acceptance decision inspectable.

  \item A repeatable robot-integration package with a robot-specific driver.

  Auto-Adapter gives integration teams a consistent package for each declared
  robot configuration. The package records the simulator assets, morphology and control
  facts, source-backed tasks, validation inputs, robot capability layer, and
  robot-specific driver used in a run. This makes the robot-specific
  integration boundary visible. It allows matched high-level controllers to
  complete tasks through the same fixed capability layer and the same fixed and
  validated robot-specific driver instead of hiding robot behaviour in one-off
  task scripts.

  \item Evidence for comparing LLM backbones and diagnosing failures.

  AutoAdapter-Bench records the LLM backbone, generation condition, robot
  configuration, outcomes for driver synthesis and capability-layer use,
  failure patterns, resource use, and per-trial evidence under matched
  conditions. These records help engineering teams compare LLM backbones,
  reproduce runs, and diagnose where an LLM-assisted robot-software workflow
  failed.

\end{enumerate}

These contributions apply only to the declared tasks, robot configurations,
generation conditions, and Direct-MuJoCo evaluation protocol. They do not
establish compliance with an industry standard, industrial certification,
real-SDK fidelity, hardware safety, sim-to-real transfer, reduced
integration time, production reliability, a universally superior LLM, or a
causal morphology effect.


\section[Thesis Structure]{Thesis Structure}
\label{sec:intro-structure}

The structure of this thesis is organised as follows:

\begin{itemize}[leftmargin=*,itemsep=0.7em]

  \item Chapter~2---Background and Literature Review. This chapter
  establishes the concepts and literature needed for the study. It defines the
  relationship between high-level controllers, robot capability layers, and
  robot-specific drivers, then reviews LLM tool use, robot-software synthesis,
  independent validation, repair, and cross-embodiment research. It concludes
  by identifying the assumption that robot capability layers, the robot-specific
  drivers that implement them, capability-level pass criteria, and executable
  validation already exist, thereby motivating the three experiments.

  \item Chapter~3---AutoAdapter-Bench: Comparing LLM Backbones in Driver
  Synthesis and Capability-Layer Use. This chapter presents Experiment~1. The
  driver-synthesis track compares LLM backbones under skeleton-assisted and
  from-scratch conditions. Within each condition, the Auto-Adapter framework,
  robot inputs, source-backed Task Library, source-task pass standards, Harness
  measurement and verdict rules, and resource budget are fixed; only
  trusted-skeleton access differs. Each backbone designs a capability layer and
  synthesises its robot-specific driver, while model-authored capability-level
  pass criteria remain evaluated outputs. The use track compares the same
  backbones using a fixed capability layer backed by the same fixed and
  validated robot-specific driver, the same high-level-controller
  implementation, and matched tasks. The two outcomes are reported separately
  and motivate Chapter~4.

  \item Chapter~4---Component Analysis of the Auto-Adapter Framework.
  This chapter presents Experiment~2. First, it tests whether an LLM can derive
  five to ten reusable capabilities from at least twenty source-backed tasks
  without a pre-authored catalogue or mapping. Second, it tests whether an LLM
  can transform reviewed task-level pass standards from traceable primary
  benchmarks or industrial-production standards into source-grounded
  capability-level pass criteria, which an implementation-blind process
  independently audits and compiles. Third, matched later runs with and without
  reviewed Experience test Continued Evolution. The evidence for each part is
  reported separately, and a generated or compilable criterion is not treated
  as reliable by itself.

  \item Chapter~5---Cross-Morphology Analysis of Robot-Specific Driver
  Synthesis for Low-Level Motion-Control Capabilities. This chapter presents
  Experiment~3. It compares the synthesis of robot-specific drivers that
  implement low-level motion-control capabilities across the declared
  morphology cohort while holding the LLM backbone, high-level-controller
  implementation, Auto-Adapter configuration, generation condition,
  evaluation protocol, and resource budget fixed. Validation outcomes, failure
  patterns, and resource use characterise morphology-associated differences.
  Because morphology co-varies with actuation and control structure, the
  chapter interprets these as associations rather than causal effects.

  \item Chapter~6---Conclusion and Future Work. The final chapter
  answers the research questions in order and synthesises the evidence from
  Chapters~3--5. It relates the findings to the scientific and bounded industry
  contributions, states the limits of the Direct-MuJoCo evaluation, and
  identifies future work arising from those limits.

\end{itemize}
